\exercise{SEIS center manifold}{4}
The SEIS model is an epidemic model in which contact between a susceptible and an infectious agent causes the susceptible agent to enter an exposed (E) state, where it carries the disease but isn't infectious yet. Exposed agents then progress to the infectious I state before eventually recovering  back to the susceptible state. The SEIS model is described by the equations 
\eqan{
\dot{E}&=&pI(1-E-I)-E \\
\dot{I}&=&E-I.
}
Our goal in this exercise is to decide whether a discontinous transition takes place at the epidemic threshold, withour computing nontrivial steady states.
\subquestion 
Find the epidemic threshold, by analyzing the stability of the trivial state.

\solution 
We consider the trivial steady state $I=E=0$ and compute the Jacobian 
\eq{
{\bf J} = \avecc{-pI-1 & p(1-E-I)-pI \\ 1 & -1}
}
and substitute the steady state, to find 
\eq{
{\bf J} = \avecc{-1 & p \\ 1 & -1}
}
Since this is two-dimensional system we can test for Hopf bifurcations with the test function ${\rm Tr}(\bf J)$, however since the trace is fixed at -2 there is no Hopf bifurcation. We check for zero eigenvalues by computing the determinant
\eq{
|{\bf J}| = 1 -p 
}
which reveals a bifurcation at $p=1$. 

\subquestion 
Compute the eigenvectors in the bifurcation point. Define new coordinates $x$ and $y$ such that $x$ is the coordinate in the direction of the eigenvector with the zero eigenvalue and $y$ is the coordinate in the direction of the other eigenvalue. (Keep $p$ in the equations for now.)

\solution 
At $p=1$ the Jacobian is 
\eq{
{\bf J} = \avecc{-1 & 1 \\ 1 & -1}
}
The eigenvectors of this matrix are 
\eq{
\vec{v_1} = \avec{1\\ 1} \qqq \vec{v_2} = \avec{1\\ -1}
}
and the corresponding eigenvalues are $\lambda_1=0$ and $\lambda_2=-2$. To translate into the new coordinate system, we write 
\eq{
\avec{E\\ I } = x \avec{1\\ 1} + y \avec{1\\ -1}
}
and therefore 
\eqa{
E &=& x+y \\
I &=& x-y. 
}
The inverse transformation is given by 
\eqa{
x &=& (E+I)/2  \\
y &=& (E-I)/2
}
and hence 
\eqa{
\dot{x} &=& (\dot{E}+\dot{I})/2 = (pI(1-E-I)-I)/2 = (p(x-y)(1-2x)-x+y)/2 \\
\dot{y} &=& (\dot{E}-\dot{I})/2 = (pI(1-E-I)-2E+I)/2 = (p(x-y)(1-x)-x-3y)/2 
}

\subquestion
Introduce a new parameter $\rho=p-1$. Use the ansatz  $y=h(x,\rho) = ax + O(2)$ to find a (very simple) approximation for the center manifold.

\solution 
We define 
\eq{\rho=p-1} 
and then substitute 
\eq{
p=\rho+1 
}
into the equations, which yields
\eqa{
\dot{x} &=& \frac{(\rho+1)(x-y)(1-2x)-x+y}{2} \\
\dot{y} &=& \frac{(\rho+1)(x-y)(1-x)-x-3y}{2}
}
Now we start with the Ansatz 
\eq{
y=ax   
}
where we ignored the higher orders for now 
If we differentiate this in time we get 
\eq{
\dot{y}=a \dot{x}   
}
We can now substitute $\dot{x}$ and $\dot{y}$ and immediately multiply by two, which yields 
\eq{
(\rho+1)(x-y)(1-x)-x-3y = a ((\rho+1)(x-y)(1-2x)-x+y)
}
We only need to match the lowest order expressions on each side which are int this case the ones that contain only one of the small variables. So as we multiply the sides out we need to keep only those terms that do not contain products of small variables, 
\eq{
x-y-x-3y = a (x-y-x+y) 
}
which simplifies to 
\eq{
-4y=0 
}
but we aren't done yet, we can substitute our Anstatz $y=ax$ again to find 
\eq{
ax=0
}
Because we are interested in small $x$ and not $x=0$, this shows $a=0$ and hence 
\eq{y=0.}

\subquestion Constrain your dynamics to the center mainfold (substitute your $y$ into $\dot{x}$). Compare the result to the normal form of the transcritical bifurcation. Is the bifurcation in the SEIS model \emph{supercritical} as in the SIS model or \emph{subcritical} as in the aSIS model with strong rewiring (see Ex.~16.11)? 

\solution 
We have a very simple approximation for $y$ indeed: $y=0$. Substituting into the equation for $\dot{x}$ yields 
\eq{
\dot{x}=\frac{(\rho+1)x(1-2x)-x}{2} = \frac{x\rho-2x^2\rho-2x^2}{2}
}
In the derivation of the transcritical normal form we showed that if there if there are nonzero $x^2$ and $\rho x$ terms then all higher order terms can be ignored as they do not change the behavior in the neighborhood of the bifurcation point. We can therefore now ignore the $x^2\rho$ term and write 
\eq{
\dot{x} = \frac{1}{2}x\rho - x^2 
}
In Ex.~16.11 we wrote the normal form of the transcritical bifurcation as 
\eq{
\dot{x} = ax + bx^2,
}
where the $a$ is not the same $a$ as our coefficient $a$ above. 
In the previous exercise we also showed that the trivial state is stable for $a<0$ and the bifurcation behaves supercritcally if $b<0$ and subcritically if $b>0$. Comparing with our SEIS system we can now see that
\eq{
a=\frac{1}{2}\rho \qqq  b= -1
}
This shows that the disease free state is stable for $\rho<0$ or equivalently $p<1$. At $p=1$ the disease free state loses stability in a transcritical bifurcation (we proved that on the center manifold, the system behaves like the transcritical normal form!). Moreover, we can say that the `supercritical' variant of this bifurcation will occur, such that the onset of the epidemic is relatively smooth, compared to the sudden explosion we saw in the aSIS model. 
Note that we showed all these things without computing the non-trivial steady state. In the SEIS model, computing the non-trivial state is quite easy, so one would not normally go this route, however for example in a Hopf bifurcation it is usually easier to distinguish subcritical and supercritical forms by computing the normal form coefficients rather than trying to compute the cycle and checking its stability.  

\subquestion 
Perhaps the linear approximation for $y$ was a bit reckless. Try again with the ansatz $y=bx^2+cx\rho + O(3)$.

\solution
Differentiating the ansatz leads to 
\eq{
\dot{y}=(2bx+c\rho)\dot{x}. 
}
By now we know that we won't need all the higher order terms. So let'S multiply out $\dot{y}$ while ignoring all terms greater than 2nd order. While we are on it note that $y$ itself is a second order term now, so we can ignore all terms where $y$ is multiplied with anything else. SO we write  
\eq{
\dot{y}= \frac{x\rho + x - y -2x^2 -x-3y}{2} + O(3) = \frac{x\rho}{2} -x^2 -2y + O(3)   
}
For $\dot{x}$ first order terms up to first order are sufficient, because the multiplication with $2bx+c\rho$ will add another order, so 
\eq{
\dot{x} = \frac{x -x}{2} + O(2) = O(2)   
}
substituting $\dot{x}$ into the equation above yields
\eq{
\dot{y} = O(3)
}
and then substituting $\dot{y}$, brings us to 
\eq{
\frac{x\rho}{2} -x^2 -2y + O(3) = O(3)
}
We cancel the higher order terms, multiply everything with two and substitute the $y=bx^2+cx\rho$ to arrive at  
\eq{
x\rho -2x^2  = 4bx^2+4cx\rho
}
By comparing the coefficients we find 
\eqa{
  1 &=& 4c \\
 -2 &=& 4b
}
and hence 
\eq{
b=-\frac{1}{2} \qqq c = \frac{1}{4}
}
such that 
\eq{
y= -\frac{1}{2}x^2 + \frac{1}{4} x \rho
}
so the center manifold is indeed curved. We can now substitute this center manifold into the equation
\eq{
\dot{x}=\frac{(\rho+1)(x-y)(1-2x)-x+y}{2} = \frac{\rho(x-y)(1-2x)-2x^2+2xy}{2}
}
Note however, that in this equation each $y$ is multiplied with at least one other small parameter, since $y=O(2)$ all terms that contain a $y$ will be at least of order 3. And hence 
\eq{
\dot{x} = \frac{x \rho -2x^2}{2} +O(3) 
}
because the transcritical normal form tells us that we can ignore terms of order three we get the same result as before
\eq{
\dot{x} = \frac{x\rho}{2} -x^2
}
In some cases the center manifold is not unique, so that different center manifolds are possible. However, these differences cannot lead to qualitatively different behavior of the bifurcations.